% Generated draft: requires mathematical and source review.
\documentclass[11pt]{article}
\usepackage[margin=25mm]{geometry}
\usepackage{fontspec}
\setmainfont{Latin Modern Roman}
\usepackage{amsmath,amssymb,mathtools,unicode-math}
\setmathfont{Latin Modern Math}
\usepackage{enumitem,needspace,xurl,hyperref}
\setlength{\parindent}{0pt}
\setlength{\parskip}{5pt}
\setlength{\emergencystretch}{3em}
\newcommand{\N}{\mathbb{N}}
\newcommand{\Z}{\mathbb{Z}}
\newcommand{\Q}{\mathbb{Q}}
\newcommand{\R}{\mathbb{R}}
\newcommand{\C}{\mathbb{C}}
\newcommand{\F}{\mathbb{F}}
\newcommand{\A}{\mathbb{A}}
\newcommand{\PP}{\mathbb{P}}
\newcommand{\E}{\mathbb{E}}
\newcommand{\eps}{\varepsilon}
\newcommand{\dd}{\,\mathrm{d}}
\begin{document}
\section*{UnsolvedMath ID: 6200030}
\textbf{Source problem code:} catalogue entry
\subsection{Definitions and mathematical statement}
Equip a finitely generated group $G$ with a Cayley graph for a finite generating set and its word metric. It is one-ended if the complement of every finite set has exactly one infinite component. A quasiline is the image of a quasi-isometric embedding of $\mathbb Z$ (or $\mathbb R$); a quasi-ray is defined similarly using $\mathbb N$ (or $[0,\infty)$). A standard coarse-separation condition for a subset $Q$ is that, for some $R$, the complement of its closed $R$-neighborhood has at least two deep components, each containing points arbitrarily far from $Q$.

Extend the following three Papasoglu results from finitely presented one-ended groups to the finitely generated setting:
\begin{enumerate}
\item Quasi-isometry invariance of the relevant JSJ decomposition.
\item A characterization of coarse separation by a quasiline in terms of a nontrivial splitting over a virtually cyclic subgroup or virtual surface-group structure.
\item Nonseparation by quasi-rays.
\end{enumerate}
A group is virtually cyclic if it has a cyclic subgroup of finite index. A splitting is a nontrivial graph-of-groups decomposition; virtual surface-group structure means a finite-index subgroup is a surface fundamental group, with the surface class still to be specified.

The source does not reproduce the three original theorem statements. In particular it does not define the JSJ category or its asserted invariance, the exact coarse-separation hypotheses and equivalence in the second result, or the hypotheses of the quasi-ray result. Its phrase ``all finitely generated groups'' also leaves unstated whether one-endedness is retained. The natural same-scope extension removes finite presentation while keeping one-endedness, but that interpretation cannot silently stand for an assertion about groups with arbitrary ends.
\subsection{Short English statement}
Do Papasoglu's JSJ, quasiline-separation, and quasi-ray-nonseparation results survive when finite presentation is dropped? The exact original statements and whether one-endedness is retained are missing from the extract.
\subsection{Sources}
\begin{itemize}
\item Misha Kapovich, \emph{Problems on Boundaries of Groups and Kleinian Groups} (2005), Problem 30, PDF page 9, attributed to Panos Papasoglu, \url{https://www.math.ucdavis.edu/~kapovich/EPR/problems.pdf}; accessed 2026-07-29, PDF-text-normalized extraction.
\item The source and research summary refer to three theorems of Papasoglu for finitely presented one-ended groups without supplying publication identifiers or their full statements.
\item Local registry \texttt{data/unsolvedmath/problems.json}, permanent UnsolvedMath ID $6200030$, source code catalogue entry, \emph{Boundaries of Groups and Kleinian Groups --- Problem 30}. Dataset attribution: Ulam AI and the UnsolvedMath Contributors.
\end{itemize}
\end{document}
